% !PRINTED none
% Abschnitt-opening page: no running head / page number visible on this page.
\begin{center}
\gesperrt{Achtzehnter Abschnitt.}

\gesperrt{Wurzeln metacyklischer Gleichungen.}

\rule{0.3\linewidth}{0.4pt}

\S.\,183.

\gesperrt{Stellung der Aufgabe. Hülfssatz.}
\end{center}

Wir haben im vorigen Abschnitte allgemeine Kennzeichen gefunden, durch die man entscheiden kann, ob eine vorgelegte Gleichung von Primzahlgrad metacyklisch ist oder nicht. Damit ist aber der Gegenstand noch bei Weitem nicht erschöpft. Es handelt sich vielmehr nach der Form der Problemstellung, wie sie von \gesperrt{Abel} herrührt, darum, ein Verfahren anzugeben, nach dem man alle metacyklischen Gleichungen finden kann. In dieser Form ist das Problem, über das \gesperrt{Abel} nur kurze Andeutungen ohne Beweise hinterlassen hat, von \gesperrt{Kronecker} aufgenommen worden, und in der Weise vollständig gelöst, dass zunächst nicht die Gleichungen selbst, sondern ihre Wurzeln gefunden werden\footnote{\gesperrt{Abel}, Sur la résolution algébrique des équations (Oeuvres complètes, ed. Sylow, tome II, p. 217), Brief an \gesperrt{Crelle} am 14. März 1826 (Oeuvres tome II, p. 266). --- \gesperrt{Kronecker}, Ueber die algebraisch auflösbaren Gleichungen. Monatsberichte der Berl. Akademie 1853, 1856. --- \gesperrt{H. Weber}, Ueber algebraisch auflösbare Gleichungen von Primzahlgrad. Sitzungsberichte der Gesellschaft zur Beförderung der Naturwissenschaften zu Marburg 1892. Die folgende Darstellung ist gegen die dort gegebene vereinfacht und in einem Punkte berichtigt.}. So hat \gesperrt{Kronecker} für jede Primzahl $n$ einen Ausdruck angegeben, der aus einem gegebenen Körper $\Omega$ durch wiederholtes Wurzelziehen abgeleitet ist, und der die doppelte
%CONT>>
